% ══════════════════════════════════════════════════════════════════════════════
\chapter*{Introduction}
\addcontentsline{toc}{chapter}{Introduction}
% ══════════════════════════════════════════════════════════════════════════════

Ce rapport présente le système de recherche d'information développé dans le cadre de
l'UV LO17 à l'UTC. Le corpus cible est constitué de 326 articles issus des bulletins
de veille technologique de l'ADIT publiés entre 2011 et 2014, fournis sous la forme
de fichiers HTML hétérogènes. Le projet consiste à construire, pas à pas sur six séances
de travaux dirigés, l'intégralité de la chaîne allant de l'extraction brute des données
jusqu'à un moteur de recherche opérationnel interrogeable en langage naturel.

La figure~\ref{fig:pipeline_global} présente l'architecture globale du système. Chaque
bloc correspond à un TD dont le livrable devient l'entrée du suivant, garantissant une
progression incrémentale et vérifiable. L'ensemble du code est écrit en Python~3.12,
géré par \texttt{uv} et structuré selon la convention \emph{src layout} de la PyPA.
La couverture de tests s'élève à 93\,\% sur 369 tests unitaires à l'issue du projet.

\begin{figure}[ht]
  \centering
  \begin{tikzpicture}[
      node distance=0.6cm and 0.3cm,
      box/.style={draw, rounded corners=3pt, fill=codegray,
                  minimum width=2.6cm, minimum height=0.85cm,
                  font=\small\ttfamily, align=center},
      arr/.style={-Stealth, thick},
    ]
    \node[box] (td1) {TD1\\corpus.xml};
    \node[box, right=of td1] (td2) {TD2\\anti-dico};
    \node[box, right=of td2] (td3) {TD3\\index inversé};
    \node[box, right=of td3] (td4) {TD4\\correcteur};
    \node[box, right=of td4] (td5) {TD5\\parser req.};
    \node[box, right=of td5] (td6) {TD6\\moteur};
    \draw[arr] (td1) -- (td2);
    \draw[arr] (td2) -- (td3);
    \draw[arr] (td3) -- (td4);
    \draw[arr] (td4) -- (td5);
    \draw[arr] (td5) -- (td6);
    \node[above=0.15cm of td1, font=\tiny\itshape] {HTML brut};
    \node[above=0.15cm of td6, font=\tiny\itshape] {requête NL};
  \end{tikzpicture}
  \caption{Architecture globale du système : chaîne de traitement en six étapes.}
  \label{fig:pipeline_global}
\end{figure}
